\documentclass[aspectratio=169,10pt]{beamer}
\usetheme{Boadilla}
\usecolortheme{dolphin}
\setbeamertemplate{navigation symbols}{}
\setbeamertemplate{footline}[frame number]
\setbeamercolor{structure}{fg=blue!55!black}
\setbeamercolor{alerted text}{fg=red!65!black}
\usepackage[T1]{fontenc}
\usepackage[utf8]{inputenc}
\usepackage{booktabs}
\usepackage{tabularx}
\usepackage{array}

\title[Autonomous Driving Benchmarking]{Autonomous Driving Benchmarking}
\subtitle{Review 1: Project Foundation and Review 2: Method and Experimental Setup}
\date{Reviews 1 and 2}

\begin{document}

\begin{frame}
  \titlepage
\end{frame}

\begin{frame}{Project overview}
\small
\begin{tabularx}{\textwidth}{>{\bfseries}p{0.22\textwidth} X}
\toprule
Project type & Reinforcement-learning project with an applied experimental comparison \\
Research area & Autonomous driving and continuous control \\
Primary method & Soft Actor-Critic (SAC) \\
Baselines & Local language-model planner and human driver \\
Formal project title & \textit{Autonomous Driving Benchmarking: A Seed-Matched Paired Evaluation of Reinforcement Learning, Large Language Models, and Human Baselines} \\
\bottomrule
\end{tabularx}
\vspace{0.7em}
The central goal is to develop and evaluate the RL policy. LLM and human driving provide comparison baselines.
\end{frame}

\begin{frame}{Motivation}
\begin{itemize}
  \item Autonomous driving requires an agent to balance lane keeping, speed, progress, and safety.
  \item Reinforcement learning can learn continuous control through repeated interaction and reward feedback.
  \item Language-model planners use a different decision process and can provide an interpretable baseline.
  \item Human control provides a practical reference for interpreting simulated-agent performance.
  \item Comparing agents on different scenarios makes it difficult to tell whether results come from the agent or scenario difficulty.
\end{itemize}
\begin{block}{Motivating idea}
Evaluate the agents in one repeatable driving task with shared scenario seeds and common episode measures.
\end{block}
\end{frame}

\begin{frame}{Problem statement}
\begin{block}{Problem}
Develop a continuous-control RL driving policy and assess its behavior relative to a local LLM planner and a human driver under matched scenario conditions.
\end{block}
\textbf{Main research question}
\begin{quote}
How does a trained SAC policy compare with a local LLM planner and a human driver on the same seeded driving scenarios?
\end{quote}
\textbf{Comparison dimensions}
\begin{itemize}
  \item Safety: collisions and road departures
  \item Driving quality: lane and speed tracking, control smoothness
  \item Task performance: route progress, completion, and time
\end{itemize}
\end{frame}

\begin{frame}{Objectives and scope}
\begin{enumerate}
  \item Formulate driving as a continuous-control reinforcement-learning task.
  \item Train a SAC policy to maintain lane position, track target speed, and make safe progress.
  \item Define LLM and human baselines that use the same driving task.
  \item Record consistent state, action, and episode telemetry.
  \item Compare agents on shared seeds using relevant measures.
\end{enumerate}
\vspace{0.5em}
\textbf{Scope:} controlled simulated driving and agent comparison. Results are specific to the task and should not be presented as evidence of real-road driving ability.
\end{frame}

\begin{frame}{Literature review: reinforcement learning for control}
\begin{block}{Soft Actor-Critic (SAC)}
Haarnoja et al. (2018) introduce an off-policy maximum-entropy RL method for continuous action spaces.
\end{block}
\begin{itemize}
  \item A replay buffer allows SAC to reuse past transitions.
  \item The entropy term encourages exploration while optimizing expected reward.
  \item Continuous actions fit steering and throttle/brake control.
  \item This project uses SAC as its primary learning method.
\end{itemize}
\vspace{0.4em}
\textbf{Design implication:} evaluate both task reward and driving measures such as safety, lane tracking, and speed error.
\end{frame}

\begin{frame}{Literature review: language models in planning}
\small
\begin{tabularx}{\textwidth}{>{\bfseries}p{0.22\textwidth} X X}
\toprule
Work & Main idea & Relevance to this project \\
\midrule
SayCan (Ahn et al., 2022) & Combines language-model planning with grounded low-level skills & Motivates separating semantic decisions from vehicle control \\
GPT-Driver (Mao et al., 2023) & Frames motion planning as language generation & Connects language models with driving decisions \\
LanguageMPC (Sha et al., 2023) & Integrates language-model decisions with model-predictive control & Motivates an LLM planner paired with a deterministic controller \\
\bottomrule
\end{tabularx}
\vspace{0.7em}
\textbf{Project baseline:} the LLM proposes a high-level maneuver; a low-level controller converts it to continuous actions. Its output and fallback status must be logged.
\end{frame}

\begin{frame}{Literature review: human baselines and fair comparison}
\begin{columns}[T]
\column{0.49\textwidth}
\textbf{Human driving}
\begin{itemize}
  \item CARLA (Dosovitskiy et al., 2017) established a widely used autonomous-driving simulation platform.
  \item A live human baseline can provide a direct reference within the same task.
  \item Human trials need consistent instructions and scenario seeds.
\end{itemize}
\column{0.49\textwidth}
\textbf{Paired evaluation}
\begin{itemize}
  \item Dem\v{s}ar (2006) discusses statistical comparisons across multiple datasets.
  \item Matching scenario seeds controls for differences in scenario difficulty.
  \item Analyze within-seed differences and report effect sizes and uncertainty.
\end{itemize}
\end{columns}
\vspace{0.4em}
These works inform the study design; they do not by themselves establish the project's novelty.
\end{frame}

\begin{frame}{Research gap and proposed contribution}
\begin{block}{Working research gap}
The project asks whether a small local LLM planner, a learned SAC policy, and a live human baseline can be compared fairly when they encounter the same seeded task and share evaluation measures.
\end{block}
\textbf{Proposed contribution}
\begin{itemize}
  \item A repeatable seeded driving task and common telemetry format
  \item SAC as the primary learned continuous-control agent
  \item LLM and human baselines evaluated under matched scenarios
  \item Per-seed analysis of safety, tracking, smoothness, and task completion
\end{itemize}
\small Treat this as a project question until the literature review has checked the cited work and related benchmarks in sufficient depth.
\end{frame}

\begin{frame}{Review 1 summary}
\begin{itemize}
  \item The project is an RL study centered on training a SAC driving policy.
  \item The comparison baselines are a local LLM planner and a human driver.
  \item Shared scenario seeds and common measures are central to a fair comparison.
  \item The literature motivates SAC, hierarchical language-model planning, human baselines, and paired evaluation.
  \item The methodology uses SAC as the main learned controller and evaluates it alongside LLM and human baselines.
\end{itemize}
\end{frame}

\begin{frame}{Review 2: system architecture}
\small
\begin{center}
\begin{tabular}{c}
\textbf{Seeded driving environment} \\
$\downarrow$ \\
\textbf{Shared observation} (speed, lane, heading, obstacle, progress, curvature) \\
$\downarrow$ \\
\textbf{Agent decision} (SAC / LLM planner / human driver) \\
$\downarrow$ \\
\textbf{Vehicle action} (steering and throttle/brake) \\
$\downarrow$ \\
\textbf{Telemetry and episode measures}
\end{tabular}
\end{center}
\begin{itemize}
  \item SAC produces continuous actions at each 20 Hz simulation step.
  \item The LLM proposes a high-level maneuver at a configured 1 Hz; a deterministic controller converts it to continuous controls.
  \item The human baseline uses keyboard control in the same task.
\end{itemize}
\end{frame}

\begin{frame}{Review 2: SAC model and learning objective}
\small
\begin{block}{Primary method}
Soft Actor-Critic is an off-policy, maximum-entropy algorithm for continuous actions. Its replay buffer reuses collected transitions during training.
\end{block}
\begin{columns}[T]
\column{0.49\textwidth}
\textbf{Observation}
\begin{itemize}
  \item Speed and target speed
  \item Lane deviation and heading error
  \item Obstacle distance and route progress
  \item Road curvature
\end{itemize}
\column{0.49\textwidth}
\textbf{Action}
\begin{itemize}
  \item Steering in $[-1,1]$
  \item Throttle/brake in $[-1,1]$
  \item Continuous control at each environment step
\end{itemize}
\end{columns}
\vspace{0.3em}
\textbf{Configured model:} MLP policy, learning rate $3\times10^{-4}$, replay buffer 50,000, batch size 256, discount factor $\gamma=0.99$.
\end{frame}

\begin{frame}{Review 2: reward design}
\small
The reward combines the task's driving objectives:
\[
R = w_p R_{progress} + w_l R_{lane} + w_v R_{speed}
  + w_s R_{safety} + w_t R_{task} + w_c R_{control}
\]
\begin{tabularx}{\textwidth}{>{\bfseries}p{0.19\textwidth} X X}
\toprule
Component & Environment signal & Intended behavior \\
\midrule
Progress & Speed and route progress & Move toward the destination \\
Lane keeping & Lane deviation and heading error & Stay near the lane center \\
Speed tracking & Error from target speed & Maintain the desired speed \\
Safety & Collision and road-departure penalties & Avoid unsafe termination \\
Task completion & Destination reward & Complete the route \\
Control & Steering penalty & Avoid excessive steering \\
\bottomrule
\end{tabularx}
\end{frame}

\begin{frame}{Review 2: experimental setup}
\begin{columns}[T]
\column{0.48\textwidth}
\textbf{Driving task}
\begin{itemize}
  \item Lightweight top-down 2D road simulator
  \item 20 Hz simulation with 40 km/h target speed
  \item Generated obstacle and route-completion signal
  \item Gymnasium environment interface
\end{itemize}
\column{0.48\textwidth}
\textbf{Scenario seeds}
\begin{itemize}
  \item Training: 20 seeds, 1000--1019
  \item Validation: 10 seeds, 1100--1109
  \item Held-out test: 20 seeds, 2000--2019
  \item Use the same held-out seeds for each agent
\end{itemize}
\end{columns}
\vspace{0.4em}
Keep test scenarios separate from model selection and tuning.
\end{frame}

\begin{frame}{Review 2: data collection and analysis workflow}
\small
\begin{tabularx}{\textwidth}{>{\bfseries}p{0.20\textwidth} X X}
\toprule
Stage & Project data or method & Review status \\
\midrule
Collection & Per-step JSON telemetry for seed, agent, state, action, termination, and LLM backend & Implemented; pilot volume is small \\
Feature extraction & Speed, lane deviation, heading error, route progress, steering, and throttle/brake & Implemented \\
Preprocessing & Validate fields and match episodes by seed before comparing agents & Full-dataset checks remain \\
EDA & Metric summaries, trajectories, speed profiles, and SAC evaluation curve & Pipeline available; pilot only \\
\bottomrule
\end{tabularx}
\vspace{0.6em}
\textbf{Episode measures:} collision, road departure, completion, route progress, lane-deviation RMSE, speed-error MAE, steering smoothness, and LLM latency/fallback status.
\end{frame}

\begin{frame}{Review 2: comparison plan and current evidence}
\begin{columns}[T]
\column{0.49\textwidth}
\textbf{Analysis plan}
\begin{itemize}
  \item Match records by scenario seed.
  \item Compare within-seed outcomes.
  \item Use paired $t$-tests or Wilcoxon tests where appropriate.
  \item Report sample counts, effect sizes, and uncertainty.
\end{itemize}
\column{0.49\textwidth}
\textbf{Pilot status}
\begin{itemize}
  \item SAC trained for 5,000 steps.
  \item Validation reward: $886.61 \pm 472.52$ over five episodes.
  \item One 80-step matched smoke pair on seed 2003 ended at the time limit without collision.
  \item Ollama was unavailable; no human run has been collected.
\end{itemize}
\end{columns}
\begin{alertblock}{Evidence limit}
The pilot is too small and variable to support conclusions about which agent performs best.
\end{alertblock}
\end{frame}

\begin{frame}{Combined review takeaways}
\begin{itemize}
  \item Review 1 defines the RL problem, objectives, literature basis, and comparison question.
  \item Review 2 specifies the SAC controller, observation and action spaces, reward, and experiment design.
  \item The study compares SAC with LLM and human baselines under matched scenario conditions.
  \item The next work is to extend training, complete baseline runs, validate the full telemetry, and analyze held-out scenarios.
\end{itemize}
\end{frame}

\begin{frame}{References}
\small
\begin{enumerate}
  \item Haarnoja, T., Zhou, A., Abbeel, P., \& Levine, S. (2018). Soft Actor-Critic: Off-Policy Maximum Entropy Deep Reinforcement Learning with a Stochastic Actor. ICML.
  \item Ahn, M. et al. (2022). Do As I Can, Not As I Say: Grounding Language in Robotic Affordances. arXiv:2204.01691.
  \item Mao, J. et al. (2023). GPT-Driver: Learning to Drive with GPT. arXiv:2310.01415.
  \item Sha, H. et al. (2023). LanguageMPC: Large Language Models as Decision Makers for Autonomous Driving. arXiv:2310.03026.
  \item Dosovitskiy, A. et al. (2017). CARLA: An Open Urban Driving Simulator. CoRL.
  \item Dem\v{s}ar, J. (2006). Statistical Comparisons of Classifiers over Multiple Data Sets. \textit{Journal of Machine Learning Research}, 7, 1--30.
\end{enumerate}
\vspace{0.3em}
Verify the bibliographic details and the project's research-gap claim against the original sources before final submission.
\end{frame}

\end{document}

